\documentclass[preview]{standalone}

\usepackage{amsmath}
\usepackage{amssymb}
\usepackage{stellar}
\usepackage{definitions}

\begin{document}

\id{probability-exercises-batch1}
\genpage

\section{Exercises}

\begin{snippetexercise}{probability-exercises-1-ex-1}{}
    We have four books, two of which are divided into \(3\) volumes and two into
    \(4\) volumes, and we place the books on a shelf.
    \begin{enumerate}
        \item In how many ways can they be placed on the shelf?
        \item In how many ways can they be placed on the shelf so that the volumes
        of each book are not separated?
    \end{enumerate}
\end{snippetexercise}

\begin{snippetsolution}{probability-exercises-1-ex-1-sol}{}
    \todo
\end{snippetsolution}

\begin{snippetexercise}{probability-exercises-1-ex-2}{}
    Given \(M\) and \(\alpha \leq M\), determine the number of solutions of the equation
    \[
        \sum_{i=1}^{\alpha} x_i = M
    \]
    where \(x_i \in \mathbb{N} \setminus \{0\}\). What changes if
    \(x_i \in \mathbb{N}\)?
\end{snippetexercise}

\begin{snippetsolution}{probability-exercises-1-ex-2-sol}{}
    \todo
\end{snippetsolution}

\begin{snippetexercise}{probability-exercises-1-ex-3}{}
    A group of \(5\) boys and \(5\) girls sits around a round table. In how many
    ways can they sit so that there are no groups of two boys sitting next to
    each other?
\end{snippetexercise}

\begin{snippetsolution}{probability-exercises-1-ex-3-sol}{}
    \todo
\end{snippetsolution}

\begin{snippetexercise}{probability-exercises-1-ex-4}{}
    On a circumference, \(8\) points are fixed. Compute:
    \begin{enumerate}
        \item how many distinct chords can be drawn between these points;
        \item how many distinct triangles can be constructed with vertices among
        the given points.
    \end{enumerate}
\end{snippetexercise}

\begin{snippetsolution}{probability-exercises-1-ex-4-sol}{}
    \todo
\end{snippetsolution}

\begin{snippetexercise}{probability-exercises-1-ex-5}{}
    An urn contains balls numbered from \(1\) to \(n\). The balls are drawn without
    replacement and arranged in order in a row. Compute the probability that the
    \(k\)-th ball is the largest ball in the urn, knowing that all the previously
    drawn balls are smaller than the \(k\)-th one.
\end{snippetexercise}

\begin{snippetsolution}{probability-exercises-1-ex-5-sol}{}
    \todo
\end{snippetsolution}

\begin{snippetexercise}{probability-exercises-1-ex-6}{}
    A basket contains \(20\) cherries, \(5\) of which are without a pit. A pig eats
    \(5\) cherries at random, and then we take one cherry at random from the basket.
    \begin{enumerate}
        \item Compute the probability that the chosen cherry has a pit.
        \item Number from \(1\) to \(5\) the cherries eaten by the pig and let the
        cherry chosen by us be the sixth cherry removed from the basket. Define the
        events
        \[
            S_i = \{\text{the \(i\)-th chosen cherry has a pit}\}
        \]
        and, given the event
        \[
            S = \bigcap_{i=1}^{6} S_i^{a_i}
        \]
        where \(a_i \in \{0,1\}\), \(i=1,\ldots,6\), and \(S_i^1 = S_i\) while
        \(S_i^0 = \widetilde{S_i}\), show that for any permutation \(\sigma\)
        of \([6]\),
        \[
            P(S) = P(S^\sigma),
        \]
        where
        \[
            S^\sigma \triangleq \bigcap_{i=1}^{6} S_{\sigma(i)}^{a_i}.
        \]
        \item Compute the probability that at least one of the cherries eaten by
        the pig contains a pit, knowing that the cherry chosen by us had a pit.
    \end{enumerate}
\end{snippetexercise}

\begin{snippetsolution}{probability-exercises-1-ex-6-sol}{}
    \todo
\end{snippetsolution}

\begin{snippetexercise}{probability-exercises-1-ex-7}{}
    An urn contains \(b\) white balls and \(r\) red balls. All the balls are drawn
    one by one. Compute the probability:
    \begin{enumerate}
        \item of drawing the first red ball exactly at the \((k+1)\)-th draw,
        where \(k \leq b+r-1\);
        \item that the last ball drawn is red.
    \end{enumerate}
\end{snippetexercise}

\begin{snippetsolution}{probability-exercises-1-ex-7-sol}{}
    \todo
\end{snippetsolution}

\begin{snippetexercise}{probability-exercises-1-ex-8}{}
    Tizio has five fair coins, of which:
    \begin{itemize}
        \item \(2\) have heads on both faces;
        \item \(1\) has tails on both faces;
        \item \(2\) are normal.
    \end{itemize}
    Tizio closes his eyes, chooses one coin at random among the five, and tosses it.
    \begin{enumerate}
        \item Compute the probability that the lower face of the tossed coin is
        heads, without opening his eyes.
        \item Tizio opens his eyes and sees that the upper face of the coin is heads.
        What is the probability that the lower face is heads?
        \item He closes his eyes again and tosses the coin. What is the probability
        that the lower face is heads?
        \item He opens his eyes again and sees that the upper face is heads. What is
        the probability that the lower face is heads?
        \item Tizio closes his eyes again, throws away the coin he tossed, chooses
        another one, and tosses it. What is the probability of getting heads?
    \end{enumerate}
\end{snippetexercise}

\begin{snippetsolution}{probability-exercises-1-ex-8-sol}{}
    \todo
\end{snippetsolution}

\begin{snippetexercise}{probability-exercises-1-ex-9}{}
    There are \(n\) urns, and the \(r\)-th urn contains \(r-1\) red balls and
    \(n-r\) black balls. An urn is chosen at random, and two draws are made without
    replacement. Let \(N_1\) and \(N_2\) denote, respectively, the events that a
    black ball is drawn in the first and in the second draw. Compute the probabilities
    of the events
    \[
        N_1
        \quad\text{and}\quad
        N_2 \mid N_1.
    \]
\end{snippetexercise}

\begin{snippetsolution}{probability-exercises-1-ex-9-sol}{}
    \todo
\end{snippetsolution}

\begin{snippetexercise}{probability-exercises-1-ex-10}{}
    Let \(E = \{1,2,\ldots,N\}\), where \(N\) is a prime number, endowed with the
    uniform distribution. Show that if \(A\) and \(B\) are two independent events,
    then at least one of \(P(A)\) and \(P(B)\) is equal to \(0\) or to \(1\).
    Equivalently, at least one of \(A\) and \(B\) is either \(\emptyset\) or
    \(\Omega\).
\end{snippetexercise}

\begin{snippetsolution}{probability-exercises-1-ex-10-sol}{}
    \todo
\end{snippetsolution}

\end{document}